%-*-tex-*-
\ifundefined{writestatus} \input status \relax \fi %
\chcode{cite}

\def\cqu{\cquote{Always verify your references.}{Martin~Joseph~Routh
(1755-1854)}
}

\chapterhead{cite}{CITATIONS\cr or REFERENCE\cr LISTS}
\intex\ provides facilities for automatic numbering and creation of {\it
citation } or {\it reference} lists. The listing format supported 
is that of the IEEE, which numbers citations sequentially in the
order that they are called out in the document. For short lists (about 40-60
citations), the actual citation forms may be in random order. For very long
lists, the citation form file must be ordered in the same manner as the
citations should appear in the document. Page referencing for citations is
not supported. 

\shead{citcomlist}{Command List}
\begintwocolumn
\ext|\citenum|
\ext|\cite|
\ext|\citeform|
\ext|\citetagvalue|
\ext|\closecitationtagfile|
\ext|\initcitetag{<integer>}|
\ext|\inputciteformfiles|
\ext|\inputcitetagfiles|
\ext|\lastcitetag{<integer>}|
\ext|\makecitationlist|
\ext|\makelongcitationlist|
\endthreecolumn

\shead{makecite}{Making Citations}
A citation is made in a document using a
format such as |[\cite{<citation tag>}]| where the |<citation tag>| may be of
any valid tag form (see Section~{\ref{validtag}}, page \pref{validtag}). 
If a |\cite| is used in internal vertical mode, for example in a section or
chapter head title, then the margin notes will disappear. If it used in any
field that may become a tag, then disaster will occur. This latter should
never happen. 
When the first |\cite| is called, the citation tag file is opened
automatically and tag generated for this citation, and the 
numerical value is placed in the text. Subsequently, a |\cite|
will 
define a new tag if it is the first time it has been called and place the
numerical value in the text. The command |\citeref{<tag>}| will place the
numerical value of the citation in the text but will neither open the
citation tag file if it is closed nor define a new tag if it is undefined. 
If a citation is to be defined but the value should not be put in the text,
then 
|{\silenttrue \cite{<citation tag>}}| should be used.


The actual citation forms are maintained  in a separate file whose
default name is |\jobname.cfm|. The command |\makecitationlist| 
does not require that the citation forms in the |.cfm| file be in the same
order and may include citation forms that are not even used. However,
because of limitations on \tex's internal memory, the maximum number of
citations in a  |.cfm| file used by |\makecitationlist| 
is limited to approximately 40 to 60. If \tex\ runs out of memory it is
necessary to use |\makelongcitationlist| instead. When using this command,
the |.cfm| file may be of any length but the citation forms must be in the
correct order, and include only those actually cited in the document. 
This must be done by some program external to \intex.


The citation form command in the |.cfm| file is 
|\citeform{<citation tag>}{<TeX citation form macro>}|. The 
|<TeX citation form macro>| is of any form but will usually be a special
|\listitem|. The actual number of the particular citation being written is
held in the command |\citetagvalue|. This command should be placed at that
place in the |<TeX citation form macro>| where the citation number should
appear. 

The first call to the |\cite| command will open the |\jobname.ctg| file where
the tags for the citations are stored. 
Then
|\initcitetag{<integer>}| where |<integer>| is one less than the lowest
number in the citation list is written. This number allows for several
independent citation lists to exist in a single document, for instance at the
end of each chapter. When the file is closed, either with a
|\makecitationlist| or a |\closecitationtagfile|, an
|\lastcitenum{<integer>}| is written. The |<integer>| is the current value of
the |\citenum|. 

To make a citation list the following is required:
\beginttverbatim
\1beginlist % citations are usually listitems ... should be group
\1frenchspacing % suppresses extra space after periods 
<initial macros \1ejectpage, make Reference header
       and set up list>
\1makecitationlist % or \1makelongcitationlist
\1endlist
\endttverbatim
The actual list of macros is inserted, in order by |\makecitationlist| or
by |\makelongcitationlist|. The latter must be used if there are too many
citations for \tex's memory. In that case the |\jobname.cfm| file must be
sorted and contain only those citations actually used. 
These commands automatically close the |\jobname.ctg| file, if it is open.

\shead{prelcite}{Prelude Citations}
Occasionally citations appear in the prelude of a document, sometimes 
through a 
section head and hence in the table of contents. 
Usually the citation numbers should be those that occur later in the text
rather than the ones that would occur naturally.
The |\cite| in the prelude causes the citation tag file to open prematurely
and the |\citenum| to be incremented incorrectly. To get around this problem,
use the following procedure:
\begintt
\begingroup % limits effect to initial section 
\inputwithcheck{<filename>.ctg} %inputs the .ctg file if it exists
\let\cite=\citeref % cites only refer now
<text with \cites>
\endgroup % where you want citations to behave normally
\endtt 


\shead{citecomforms}{Command Forms}
\beginblockmode
\ext\@|\citenum|
\nbr
|\citenum| is the counter that holds the current largest citation number
used. This number is equal to or larger than that set by
|\initcitetag{<integer>}|. 

\mbr
\ext\@|\cite{<citation tag>}|
\nbr
This command places the citation number, corresponding to the 
|<citation tag>| in the text, without any surrounding brackets or braces.
If the |<citation tag>| has not been defined, or if it was defined for a
previous citation list, the |\citenum| is increased by one, and a tag is
created with that new value. Multiple citations may be called in a single
place by  |[\cite{<citetag1>}, \cite{<citetag2>}, ...]|. However, there is
{\it no} sorting capability so that the numbers may appear out of sequential
order. The first |\cite| opens the citation tag file. 
|{\silenttrue \cite{<citetag>}}| will define and increment the citation
counter but will not insert the citation value in the text. 

\mbr
\ext\@|\citeform{<citation tag>}{<TeX macro form>}|
\nbr
This is the form that must be used in the file |\jobname.cfm| 
file that actually contains the list of citations. The |<citation tag>| is
the same on used in the text when the citation is called out. The 
|<TeX macro form>| is the \tex\ macro that will actually be inserted into the
reference list. This must contain all the information needed to produce the
final form of the citation. Since the actual citation number is being
supplied by \intex, it is necessary to include the \tex\ command
|\citetagvalue| in the place required in the text. For instance, if the
|\listitem| command was being used as the basic form, and if the citation tag
were |[Knut84]|, then an appropriate |\citeform| for the \texbook\ would be
\begintt
\citeform{[Knut84]}{\listitem \citetagvalue D.E. Knuth, 
        {\bf \TeX book}, {\it Addison Wesley Pub. Co.}, 1984}
\endtt


\mbr
\ext\@|\citetagvalue|
\nbr
This command contains, at the time of a |\citeform| execution, the number of
that particular citation.


\mbr
\ext\@|\closecitationtagfile|
\nbr
This closes the |\jobname.ctg| file, if it is open. Since the |.ctg|
file is closed
automatically whenever |\makecitationlist| is called, this command is seldom
used. 

\mbr
\ext\@|\initcitetag{<integer>}|
\nbr
This command is written into the |.ctg| file when it is opened. It
establishes the minimum value of |\citenum| for this list. The actual
reference value is determined by subtracting this value from the
|<tagvalue>|.



\mbr
\ext\@|\inputciteformfiles 
          [=] { \inputwithcheck{<first form file>}
                  ...
                \inputwithcheck{<last form  file>} }|
\nbr
This is a token list of files, each prefaced by 
|\inputwithcheck|,  that contain the citations for a particular citation
list. The command |\inputwitcheck| determines whether the file exists
before trying to open it. The default is |\inputwithcheck{\jobname.cfm}|.



\ext\@|\inputcitetagfiles 
          [=] { \inputwithcheck{<first tag file>}
                  ...
                \inputwithcheck{<last tag file>} }|
\nbr
This is a token list of files, each prefaced by 
|\inputwithcheck|,  that contain the citation tags for a particular citation
list. The command |\inputwitcheck| determines whether the file exists
before trying to open it. The default is |\inputwithcheck{\jobname.ctg}|.
The |\jobname.ctg| file contains the tags for the citation list.
These tags are in order of call out and numbered
sequentially. It also includes the |\initcitetag| command. 

\mbr
\ext\@|\lastcitetag{<integer>}|
\nbr
This command is written into the |.ctg| file when it is 
closed either by a |\closecitationtagfile| or a 
|\makecitationlist|. It is used to reset |\citenum| when using successive 
independent 
citation lists.



\mbr
\ext\@|\makecitationlist|
\nbr
This command reads in the |\jobname.cfm| file. It then reads the
|\jobname.ctg| file and converts the list of citation tags into 
a correctly ordered list of citations. The citation forms  in the 
|\jobname.cfm| file may be in any order and may include ones that are not
used. 
Any  formatting control commands such as 
|\beginlist ... \endlist| must be placed around it.


\mbr
\ext\@|\makelongcitationlist|
\nbr
This command reads in the |\jobname.cfm| file and assigns citation numbers 
to the citation forms in the order that they appear. Errors from those
indicated by the |\citetagvalue| are reported. 
 The citation forms  in the 
|\jobname.cfm| file {\bf must be in the same order as called out in the
document and include only those that are actually used.}
Any  formatting control commands such as 
|\beginlist ... \endlist| must be placed around it.


\endblockmode
\ejectpage

\done